% Setting up the document and importing necessary packages
\documentclass[12pt,a4paper]{article}

% Chinese support packages
\usepackage{xeCJK}
\usepackage{fontspec}
\setCJKmainfont{Noto Serif CJK TC}
\setCJKsansfont{Noto Sans CJK TC}
\setCJKmonofont{Noto Sans CJK TC}

% Other necessary packages
\usepackage{geometry}
\usepackage{titlesec}
\usepackage{enumitem}
\usepackage{color}
\usepackage{colortbl}
\usepackage{tcolorbox}
\usepackage{booktabs}
\usepackage{longtable}
\usepackage{array}
\usepackage{graphicx}
\usepackage{tikz}
\usepackage{mdframed}
\usepackage{fancyhdr}
\usepackage{hyperref}
\usepackage{multicol}
\usepackage{datetime2}
\usepackage{natbib}
\usepackage{rotating}
\usepackage{adjustbox}

\hypersetup{
    colorlinks=true,
    linkcolor=emergency_blue,
    citecolor=emergency_blue,
    urlcolor=emergency_blue,
    pdfborder={0 0 0}
}

% Page setup
\geometry{
    top=2.5cm,
    bottom=2.5cm,
    left=2cm,
    right=2cm,
    headheight=15pt
}

% Color definitions
\definecolor{emergency_red}{RGB}{186,24,27}
\definecolor{emergency_blue}{RGB}{0,114,188}
\definecolor{emergency_gray}{RGB}{120,120,120}
\definecolor{highlight_yellow}{RGB}{255,250,205}
\definecolor{light_blue}{RGB}{173,216,230}
\definecolor{light_green}{RGB}{144,238,144}
\definecolor{light_orange}{RGB}{255,218,185}

% Title format settings
\titleformat{\section}[block]{\Large\bfseries\color{emergency_red}}{\thesection}{10pt}{}
\titleformat{\subsection}[block]{\large\bfseries\color{emergency_blue}}{\thesubsection}{8pt}{}
\titleformat{\subsubsection}[block]{\normalsize\bfseries\color{emergency_gray}}{\thesubsubsection}{6pt}{}

% Custom environment
\newmdenv[
    linecolor=emergency_red,
    backgroundcolor=highlight_yellow,
    linewidth=2pt,
    topline=true,
    bottomline=true,
    leftline=true,
    rightline=true,
    innertopmargin=10pt,
    innerbottommargin=10pt,
    innerleftmargin=10pt,
    innerrightmargin=10pt
]{importantbox}

\newcommand{\note}[1]{
    \par
    \noindent
    \begin{tcolorbox}[
        colback=emergency_blue!5!white,
        colframe=emergency_blue,
        title=\textbf{重點提醒},
        fonttitle=\bfseries,
        boxrule=0.5pt,
        arc=3pt,
        left=6pt,
        right=6pt,
        top=6pt,
        bottom=6pt
    ]
    #1
    \end{tcolorbox}
}

% Header and footer settings
\pagestyle{fancy}
\fancyhf{}
\fancyhead[L]{\textcolor{emergency_red}{心血管中心教學計畫}}
\fancyhead[R]{\textcolor{emergency_blue}{Clinical  Program}}
\fancyfoot[C]{\textcolor{emergency_gray}{\thepage}}
\renewcommand{\headrulewidth}{1pt}
\renewcommand{\headrule}{\hbox to\headwidth{\color{emergency_red}\leaders\hrule height \headrulewidth\hfill}}

% Date format
\DTMsetstyle{yyyy-mm-dd}
\newcommand{\mydate}{\DTMtoday}

% Document start
\begin{document}

% Title page
\begin{titlepage}
    \centering
    {\LARGE\textcolor{emergency_red}{\textbf{心血管中心教學計畫}}\par}
    \vspace{1cm}
    {\huge\bfseries 見習醫學生心臟內科臨床跟隨計畫\par}
    \vspace{0.5cm}
    {\Large\textcolor{emergency_blue}{\textit{Clinical Shadowing Program in Cardiology}}\par}
    \vspace{1cm}
    
    \begin{tcolorbox}[
        colback=emergency_blue!5!white,
        colframe=emergency_blue,
        width=0.8\textwidth,
        arc=10pt,
        boxrule=1pt
    ]
    \centering
    {\large\textbf{計畫期間}\par}
    \vspace{0.3cm}
    2025年11月17日 - 2025年12月5日\\
    （共三週完整臨床見習）\\
    \vspace{0.5cm}
    {\large\textbf{指導醫師}\par}
    \vspace{0.3cm}
    洪國軒醫師、蕭喻中醫師\\
    心臟內科主治醫師\\
    \vspace{0.5cm}
    {\large\textbf{見習學生}\par}
    \vspace{0.3cm}
    黃牧凡\\
    （醫學系六年級學生）
    \end{tcolorbox}

    \vfill
    
    {\large 心血管中心\\ 發布日期：\mydate}
    
    \vspace{0.5cm}
\end{titlepage}

\newpage
\tableofcontents
\newpage

\section{計畫概述}

\subsection{計畫目標}
本臨床跟隨計畫旨在讓見習醫學生深入了解心臟內科的臨床實務，透過直接參與臨床工作，獲得寶貴的實際經驗。

\textbf{主要學習目標：}
\begin{itemize}[nosep]
    \item 熟悉心臟內科門診流程和病患評估
    \item 學習心導管手術的基本操作和觀察
    \item 掌握心血管超音波檢查技術
    \item 了解住院病患的照護流程
    \item 培養臨床思維和醫病溝通技巧
    \item 學習電子病歷系統操作
\end{itemize}

\subsection{計畫特色}
參考國外頂尖醫學中心（Mayo Clinic、Cleveland Clinic、Johns Hopkins）的clinical shadowing模式，結合本院實際臨床環境設計：

\begin{itemize}[nosep]
    \item \textbf{全方位體驗}：涵蓋門診、病房、導管室、超音波室
    \item \textbf{實作學習}：親自操作超音波、參與穿刺技術
    \item \textbf{急診待命}：參與STEMI緊急處置
    \item \textbf{系統學習}：學習醫院資訊系統操作
    \item \textbf{個人指導}：一對一mentoring模式
\end{itemize}

\begin{importantbox}
\textbf{重要提醒：}本計畫要求學生具備基本醫學知識，並已完成基礎臨床技能訓練。學生需嚴格遵守醫院感染控制規範和患者隱私保護政策。
\end{importantbox}

\section{預習準備事項}

\subsection{必要預習內容}
開始臨床跟隨前，學生需完成以下預習：

\subsubsection{Portal電腦系統操作}
\begin{itemize}[nosep]
    \item 電子病歷（EMR）基本操作
    \item 醫囑輸入系統（CPOE）使用方法
    \item 檢查報告查詢系統
    \item 影像檢視系統（PACS）
    \item 藥物處方系統
    \item 實驗室結果查詢
\end{itemize}

\subsubsection{心導管影像判讀基礎}
\begin{itemize}[nosep]
    \item 冠狀動脈解剖學
    \item 基本造影角度認識
    \item 常見病灶判讀
    \item 血管內超音波（IVUS）基礎
    \item 光學同調斷層掃描（OCT）介紹
\end{itemize}

\subsubsection{臨床基礎知識}
\begin{itemize}[nosep]
    \item 常見心血管疾病診斷標準
    \item 心電圖判讀基礎
    \item 實驗室數據解讀
    \item 心血管藥物分類和作用機轉
    \item 心臟衰竭分級和治療指引
\end{itemize}

\vspace{0.5cm}
\note{
建議學生在計畫開始前一週完成相關預習，並準備筆記本記錄學習心得。
}

\section{第一週詳細時程表}

\subsection{週一（11月17日）- 門診見習 + 病房迴診}

\begin{longtable}{|>{\raggedright\arraybackslash}p{2cm}|>{\raggedright\arraybackslash}p{4cm}|>{\raggedright\arraybackslash}p{4cm}|>{\raggedright\arraybackslash}p{4cm}|}
\hline
\rowcolor{light_blue}
\textbf{時間} & \textbf{活動內容} & \textbf{學習重點} & \textbf{評估方式} \\
\hline
08:00-08:30 & 
\begin{itemize}[nosep]
\item 報到和orientation
\item Portal系統登入練習
\item 門診流程介紹
\end{itemize} & 
\begin{itemize}[nosep]
\item 醫院規範了解
\item 系統操作熟悉
\item 專業態度建立
\end{itemize} & 
口頭問答 \\
\hline
08:30-12:00 & 
\begin{itemize}[nosep]
\item 跟隨門診看診
\item 觀察病史詢問技巧
\item 學習理學檢查
\item 協助輸入病歷
\end{itemize} & 
\begin{itemize}[nosep]
\item 醫病溝通技巧
\item 症狀評估方法
\item 診斷思維過程
\item EMR系統使用
\end{itemize} & 
實際操作觀察 \\
\hline
12:00-13:00 & 午餐時間 & 與指導醫師討論上午案例 & 案例討論 \\
\hline
13:00-15:00 & 
\begin{itemize}[nosep]
\item 繼續門診見習
\item 學習診斷決策
\item 練習開立檢查
\end{itemize} & 
\begin{itemize}[nosep]
\item 鑑別診斷思考
\item 檢查適應症
\item 治療計畫制定
\end{itemize} & 
病例分析 \\
\hline
15:00-18:00 & 
\begin{itemize}[nosep]
\item 病房迴診
\item 學習床邊評估
\item 觀察住院醫療流程
\item 參與team meeting
\end{itemize} & 
\begin{itemize}[nosep]
\item 住院病患照護
\item 團隊合作模式
\item 病情追蹤方法
\item 多專科會診
\end{itemize} & 
團隊討論參與 \\
\hline
18:00-22:00 & 
\begin{itemize}[nosep]
\item 晚餐休息
\item 準備值班
\item 複習白天學習內容
\end{itemize} & 
\begin{itemize}[nosep]
\item 自主學習能力
\item 知識整合
\item 準備急診待命
\end{itemize} & 
學習心得記錄 \\
\hline
22:00-08:00 & 
\textcolor{emergency_red}{\textbf{STEMI值班待命}} & 
\begin{itemize}[nosep]
\item 急診處置流程
\item 緊急導管手術
\item 危重病患照護
\end{itemize} & 
急診案例參與 \\
\hline
\end{longtable}

\subsection{週二（11月18日）- 導管室實習 + 病房接診}

\begin{longtable}{|>{\raggedright\arraybackslash}p{2cm}|>{\raggedright\arraybackslash}p{4cm}|>{\raggedright\arraybackslash}p{4cm}|>{\raggedright\arraybackslash}p{4cm}|}
\hline
\rowcolor{light_green}
\textbf{時間} & \textbf{活動內容} & \textbf{學習重點} & \textbf{評估方式} \\
\hline
08:00-08:30 & 
\begin{itemize}[nosep]
\item 導管室orientation
\item 感染控制訓練
\item 輻射防護講解
\end{itemize} & 
\begin{itemize}[nosep]
\item 無菌技術概念
\item 安全防護措施
\item 導管室規範
\end{itemize} & 
安全規範測驗 \\
\hline
08:30-12:00 & 
\begin{itemize}[nosep]
\item 上台刷手練習
\item 觀察動脈穿刺技術
\item 學習導管操作
\item 參與冠狀動脈造影
\end{itemize} & 
\begin{itemize}[nosep]
\item 無菌操作技術
\item 血管穿刺方法
\item 導管選擇原則
\item 造影技術應用
\end{itemize} & 
實際操作評分 \\
\hline
12:00-13:00 & 午餐時間 & 討論導管案例 & 案例分析 \\
\hline
13:00-15:00 & 
\begin{itemize}[nosep]
\item 繼續導管手術觀摩
\item 學習支架置放
\item 觀察併發症處理
\end{itemize} & 
\begin{itemize}[nosep]
\item 介入治療技術
\item 器材選擇考量
\item 風險評估能力
\end{itemize} & 
技術理解測驗 \\
\hline
15:00-18:00 & 
\begin{itemize}[nosep]
\item 病房接新病人
\item 學習病史收集
\item 練習理學檢查
\item 協助安排檢查
\end{itemize} & 
\begin{itemize}[nosep]
\item 入院評估流程
\item 檢查優先順序
\item 初步診斷能力
\item 醫療團隊溝通
\end{itemize} & 
病史收集評估 \\
\hline
22:00-08:00 & 
\textcolor{emergency_red}{\textbf{STEMI值班待命}} & 
急診心導管處置 & 
緊急處置參與 \\
\hline
\end{longtable}

\subsection{週三（11月19日）- 門診見習 + 病房迴診}

\begin{longtable}{|>{\raggedright\arraybackslash}p{2cm}|>{\raggedright\arraybackslash}p{4cm}|>{\raggedright\arraybackslash}p{4cm}|>{\raggedright\arraybackslash}p{4cm}|}
\hline
\rowcolor{light_orange}
\textbf{時間} & \textbf{活動內容} & \textbf{學習重點} & \textbf{評估方式} \\
\hline
08:00-12:00 & 
\begin{itemize}[nosep]
\item 獨立協助門診
\item 預先評估病患
\item 準備病歷摘要
\item 學習臨床決策
\end{itemize} & 
\begin{itemize}[nosep]
\item 獨立思考能力
\item 時間管理技巧
\item 優先順序判斷
\item 效率提升方法
\end{itemize} & 
獨立作業評估 \\
\hline
12:00-13:00 & 午餐時間 & 複習前兩天學習內容 & 綜合討論 \\
\hline
13:00-15:00 & 
\begin{itemize}[nosep]
\item 慢性病追蹤門診
\item 學習長期照護
\item 觀察藥物調整
\end{itemize} & 
\begin{itemize}[nosep]
\item 慢性病管理
\item 藥物治療監測
\item 生活品質評估
\end{itemize} & 
慢病管理討論 \\
\hline
15:00-18:00 & 
\begin{itemize}[nosep]
\item 專責病房迴診
\item 學習病程記錄
\item 參與出院規劃
\item 家屬衛教觀摩
\end{itemize} & 
\begin{itemize}[nosep]
\item 病程追蹤技巧
\item 出院準備流程
\item 衛教溝通方法
\item 後續追蹤安排
\end{itemize} & 
病程記錄練習 \\
\hline
22:00-08:00 & 
\textcolor{emergency_red}{\textbf{STEMI值班待命}} & 
急診處置流程 & 
急診能力評估 \\
\hline
\end{longtable}

\subsection{週四（11月20日）- 導管室進階實習}

\begin{longtable}{|>{\raggedright\arraybackslash}p{2cm}|>{\raggedright\arraybackslash}p{4cm}|>{\raggedright\arraybackslash}p{4cm}|>{\raggedright\arraybackslash}p{4cm}|}
\hline
\rowcolor{light_blue}
\textbf{時間} & \textbf{活動內容} & \textbf{學習重點} & \textbf{評估方式} \\
\hline
08:00-12:00 & 
\begin{itemize}[nosep]
\item 複雜案例觀摩
\item 學習靜脈穿刺
\item 右心導管檢查
\item 暫時性節律器置放
\end{itemize} & 
\begin{itemize}[nosep]
\item 高難度技術觀摩
\item 靜脈穿刺技巧
\item 血流動力學評估
\item 緊急處置能力
\end{itemize} & 
技術操作評估 \\
\hline
12:00-13:00 & 午餐時間 & 導管技術討論 & 技術問答 \\
\hline
13:00-17:00 & 
\begin{itemize}[nosep]
\item 週四下午門診
\item 導管術後追蹤
\item 學習併發症評估
\item 討論治療計畫
\end{itemize} & 
\begin{itemize}[nosep]
\item 術後照護流程
\item 併發症識別
\item 長期追蹤規劃
\item 生活型態諮詢
\end{itemize} & 
術後評估練習 \\
\hline
17:00-18:00 & 
\begin{itemize}[nosep]
\item 第一週總結
\item 學習心得分享
\item 問題討論解答
\end{itemize} & 
\begin{itemize}[nosep]
\item 自我反思能力
\item 學習成效評估
\item 改進方向確認
\end{itemize} & 
週總結報告 \\
\hline
22:00-08:00 & 
\textcolor{emergency_red}{\textbf{STEMI值班待命}} & 
綜合急診能力 & 
急診處置評估 \\
\hline
\end{longtable}

\subsection{週五（11月21日）- 超音波實習 + 床邊超音波}

\begin{longtable}{|>{\raggedright\arraybackslash}p{2cm}|>{\raggedright\arraybackslash}p{4cm}|>{\raggedright\arraybackslash}p{4cm}|>{\raggedright\arraybackslash}p{4cm}|}
\hline
\rowcolor{light_green}
\textbf{時間} & \textbf{活動內容} & \textbf{學習重點} & \textbf{評估方式} \\
\hline
08:00-12:00 & 
\begin{itemize}[nosep]
\item 心血管超音波室
\item 學習Echo標準切面
\item 練習Apical 4C view
\item 練習PLAX view
\item 觀摩完整檢查
\end{itemize} & 
\begin{itemize}[nosep]
\item 超音波操作技術
\item 標準切面獲取
\item 影像品質控制
\item 基本測量方法
\item 異常發現識別
\end{itemize} & 
實際操作評分 \\
\hline
12:00-13:00 & 午餐時間 & 超音波影像判讀討論 & 影像判讀測驗 \\
\hline
13:00-15:00 & 
\begin{itemize}[nosep]
\item 週邊血管超音波
\item 學習PAOD評估
\item 學習DVT檢查
\item 動靜脈血流評估
\end{itemize} & 
\begin{itemize}[nosep]
\item 血管超音波技術
\item 都卜勒應用
\item 病理診斷標準
\item 血流評估方法
\end{itemize} & 
血管超音波操作 \\
\hline
15:00-18:00 & 
\begin{itemize}[nosep]
\item 病房床邊超音波
\item 學習便攜式設備
\item 緊急超音波應用
\item 實際病患檢查
\end{itemize} & 
\begin{itemize}[nosep]
\item 床邊診斷技術
\item 急診超音波應用
\item 快速評估能力
\item 臨床整合思維
\end{itemize} & 
床邊技能評估 \\
\hline
22:00-08:00 & 
\textcolor{emergency_red}{\textbf{STEMI值班待命}} & 
急診超音波應用 & 
緊急診斷能力 \\
\hline
\end{longtable}

\section{第二週詳細時程表}

第二週將重複第一週的模式，但增加更多獨立操作機會和進階技能學習：

\begin{importantbox}
\textbf{第二週重點：}
\begin{itemize}[nosep]
\item 增加學生獨立操作時間
\item 進階導管技術觀摩
\item 複雜病例討論
\item 研究方法介紹
\item 專題報告準備
\item 職涯規劃討論
\end{itemize}
\end{importantbox}

\subsection{進階學習目標}
\begin{itemize}[nosep]
    \item 獨立完成基本echo檢查
    \item 協助簡單導管操作
    \item 參與疑難病例討論
    \item 學習研究方法和統計分析
    \item 準備案例報告presentation
    \item 討論心臟科專科訓練路徑
\end{itemize}

\section{第三週詳細時程表}

第三週為綜合應用與總結階段，著重於整合學習成果和準備最終評估：

\begin{importantbox}
\textbf{第三週重點：}
\begin{itemize}[nosep]
\item 整合前兩週所學技能
\item 準備病例報告presentation
\item 完成技能評估
\item 參與綜合臨床討論
\item 最終評估與回饋
\item 未來發展規劃討論
\end{itemize}
\end{importantbox}

\subsection{週一（12月1日）- 整合應用與案例準備}

\begin{longtable}{|>{\raggedright\arraybackslash}p{2cm}|>{\raggedright\arraybackslash}p{4cm}|>{\raggedright\arraybackslash}p{4cm}|>{\raggedright\arraybackslash}p{4cm}|}
\hline
\rowcolor{light_blue}
\textbf{時間} & \textbf{活動內容} & \textbf{學習重點} & \textbf{評估方式} \\
\hline
08:00-12:00 & 
\begin{itemize}[nosep]
\item 獨立門診病患評估
\item 病例報告準備
\item 文獻查找和整理
\end{itemize} & 
\begin{itemize}[nosep]
\item 臨床決策能力
\item 文獻整合能力
\item 報告準備技巧
\end{itemize} & 
病例分析評估 \\
\hline
13:00-18:00 & 
\begin{itemize}[nosep]
\item 病房巡迴照護
\item 疑難病例討論
\item 指導醫師諮詢
\end{itemize} & 
\begin{itemize}[nosep]
\item 整合照護能力
\item 跨專業合作
\item 臨床思維深化
\end{itemize} & 
團隊討論參與 \\
\hline
\end{longtable}

\subsection{週二（12月2日）- 進階技能評估}

\begin{longtable}{|>{\raggedright\arraybackslash}p{2cm}|>{\raggedright\arraybackslash}p{4cm}|>{\raggedright\arraybackslash}p{4cm}|>{\raggedright\arraybackslash}p{4cm}|}
\hline
\rowcolor{light_green}
\textbf{時間} & \textbf{活動內容} & \textbf{學習重點} & \textbf{評估方式} \\
\hline
08:00-12:00 & 
\begin{itemize}[nosep]
\item 超音波技能評估
\item 獨立操作測試
\item 影像判讀考核
\end{itemize} & 
\begin{itemize}[nosep]
\item 超音波熟練度
\item 影像品質評估
\item 診斷能力測試
\end{itemize} & 
實際操作評分 \\
\hline
13:00-18:00 & 
\begin{itemize}[nosep]
\item 導管室技能複習
\item 進階案例討論
\item 併發症處理學習
\end{itemize} & 
\begin{itemize}[nosep]
\item 導管技能整合
\item 危機處理能力
\item 團隊協作評估
\end{itemize} & 
綜合能力評估 \\
\hline
\end{longtable}

\subsection{週三（12月3日）- 研究與學術能力}

\begin{longtable}{|>{\raggedright\arraybackslash}p{2cm}|>{\raggedright\arraybackslash}p{4cm}|>{\raggedright\arraybackslash}p{4cm}|>{\raggedright\arraybackslash}p{4cm}|}
\hline
\rowcolor{light_orange}
\textbf{時間} & \textbf{活動內容} & \textbf{學習重點} & \textbf{評估方式} \\
\hline
08:00-12:00 & 
\begin{itemize}[nosep]
\item 研究方法討論
\item 統計分析介紹
\item 論文寫作指導
\end{itemize} & 
\begin{itemize}[nosep]
\item 研究設計概念
\item 數據分析能力
\item 學術寫作技巧
\end{itemize} & 
討論參與評估 \\
\hline
13:00-18:00 & 
\begin{itemize}[nosep]
\item Journal club參與
\item 病例報告討論
\item Presentation練習
\end{itemize} & 
\begin{itemize}[nosep]
\item 批判性思考
\item 簡報技巧
\item 溝通表達能力
\end{itemize} & 
口頭報告評估 \\
\hline
\end{longtable}

\subsection{週四（12月4日）- 病例報告與最終評估}

\begin{longtable}{|>{\raggedright\arraybackslash}p{2cm}|>{\raggedright\arraybackslash}p{4cm}|>{\raggedright\arraybackslash}p{4cm}|>{\raggedright\arraybackslash}p{4cm}|}
\hline
\rowcolor{light_blue}
\textbf{時間} & \textbf{活動內容} & \textbf{學習重點} & \textbf{評估方式} \\
\hline
08:00-10:00 & 
\begin{itemize}[nosep]
\item 最終報告準備
\item 簡報資料整理
\item 預演練習
\end{itemize} & 
\begin{itemize}[nosep]
\item 資料組織能力
\item 時間管理
\item 壓力調適
\end{itemize} & 
準備度評估 \\
\hline
10:00-12:00 & 
\begin{itemize}[nosep]
\item 正式病例報告
\item Q\&A互動
\item 指導醫師回饋
\end{itemize} & 
\begin{itemize}[nosep]
\item 簡報表現
\item 問題回應能力
\item 專業態度
\end{itemize} & 
口頭報告評分 \\
\hline
13:00-18:00 & 
\begin{itemize}[nosep]
\item 綜合能力測驗
\item 360度評估收集
\item 學習成果檢討
\end{itemize} & 
\begin{itemize}[nosep]
\item 知識整合程度
\item 多方面能力評估
\item 學習反思能力
\end{itemize} & 
綜合評估 \\
\hline
\end{longtable}

\subsection{週五（12月5日）- 總結與展望}

\begin{longtable}{|>{\raggedright\arraybackslash}p{2cm}|>{\raggedright\arraybackslash}p{4cm}|>{\raggedright\arraybackslash}p{4cm}|>{\raggedright\arraybackslash}p{4cm}|}
\hline
\rowcolor{light_green}
\textbf{時間} & \textbf{活動內容} & \textbf{學習重點} & \textbf{評估方式} \\
\hline
08:00-10:00 & 
\begin{itemize}[nosep]
\item 學習成果總結
\item 證書頒發準備
\item 推薦信討論
\end{itemize} & 
\begin{itemize}[nosep]
\item 自我評估能力
\item 未來規劃思考
\item 專業發展方向
\end{itemize} & 
最終回饋 \\
\hline
10:00-12:00 & 
\begin{itemize}[nosep]
\item 職涯發展諮詢
\item 住院醫師申請指導
\item 研究機會討論
\end{itemize} & 
\begin{itemize}[nosep]
\item 生涯規劃
\item 申請策略
\item 個人優勢分析
\end{itemize} & 
諮詢討論 \\
\hline
13:00-15:00 & 
\begin{itemize}[nosep]
\item 結業典禮
\item 心得分享
\item 感謝與道別
\end{itemize} & 
\begin{itemize}[nosep]
\item 學習經驗整理
\item 人際關係建立
\item 專業態度展現
\end{itemize} & 
綜合表現 \\
\hline
\end{longtable}

\section{學習評估方式}

\subsection{日常評估}
\begin{itemize}[nosep]
    \item \textbf{直接觀察評估（Direct Observation）}：每日實際操作評分
    \item \textbf{案例基礎討論（Case-based Discussion）}：病例分析能力
    \item \textbf{多來源回饋（Multi-source Feedback）}：團隊成員評價
    \item \textbf{學習歷程記錄（Portfolio）}：學習心得和反思
\end{itemize}

\subsection{評估工具}
\begin{longtable}{|>{\raggedright\arraybackslash}p{3cm}|>{\raggedright\arraybackslash}p{4cm}|>{\raggedright\arraybackslash}p{4cm}|>{\raggedright\arraybackslash}p{3cm}|}
\hline
\rowcolor{light_blue}
\textbf{評估項目} & \textbf{評估方法} & \textbf{評估標準} & \textbf{權重} \\
\hline
臨床技能 & 
\begin{itemize}[nosep]
\item 超音波操作
\item 導管技術觀摩
\item 穿刺技術練習
\end{itemize} & 
\begin{itemize}[nosep]
\item 操作熟練度
\item 安全意識
\item 學習態度
\end{itemize} & 
30\% \\
\hline
知識應用 & 
\begin{itemize}[nosep]
\item 病例討論
\item 診斷推理
\item 治療計畫
\end{itemize} & 
\begin{itemize}[nosep]
\item 知識整合能力
\item 臨床思維邏輯
\item 決策合理性
\end{itemize} & 
25\% \\
\hline
溝通技巧 & 
\begin{itemize}[nosep]
\item 醫病溝通
\item 團隊合作
\item 家屬衛教
\end{itemize} & 
\begin{itemize}[nosep]
\item 溝通效果
\item 同理心表現
\item 專業態度
\end{itemize} & 
20\% \\
\hline
學習態度 & 
\begin{itemize}[nosep]
\item 主動學習
\item 問題提出
\item 反思能力
\end{itemize} & 
\begin{itemize}[nosep]
\item 學習積極性
\item 自我改進
\item 成長軌跡
\end{itemize} & 
15\% \\
\hline
系統操作 & 
\begin{itemize}[nosep]
\item EMR使用
\item PACS操作
\item 資料查詢
\end{itemize} & 
\begin{itemize}[nosep]
\item 操作熟練度
\item 效率表現
\item 準確性
\end{itemize} & 
10\% \\
\hline
\end{longtable}

\section{安全規範和注意事項}

\subsection{感染控制}
\begin{itemize}[nosep]
    \item 嚴格遵守手部衛生規範
    \item 正確使用個人防護設備
    \item 了解隔離預防措施
    \item 遵守無菌技術原則
\end{itemize}

\subsection{輻射防護}
\begin{itemize}[nosep]
    \item 導管室內必須穿著鉛衣
    \item 佩戴輻射劑量監測章
    \item 保持安全距離
    \item 了解輻射防護三原則
\end{itemize}

\subsection{患者隱私保護}
\begin{itemize}[nosep]
    \item 嚴格遵守醫療隱私法規
    \item 不得洩露患者資訊
    \item 適當的討論場所選擇
    \item 尊重患者權益
\end{itemize}

\subsection{緊急應變}
\begin{itemize}[nosep]
    \item 熟悉緊急聯絡方式
    \item 了解急救設備位置
    \item 知曉緊急撤離路線
    \item 掌握基本急救技能
\end{itemize}

\vspace{0.5cm}
\note{
學生在任何情況下都不得單獨進行侵入性操作，必須在指導醫師直接監督下學習。
}

\section{技能學習檢核表}

\subsection{超音波技能檢核}

\begin{longtable}{|>{\raggedright\arraybackslash}p{3.5cm}|>{\raggedright\arraybackslash}p{1.5cm}|>{\raggedright\arraybackslash}p{1.5cm}|>{\raggedright\arraybackslash}p{1.5cm}|>{\raggedright\arraybackslash}p{1.5cm}|>{\raggedright\arraybackslash}p{4cm}|}
\hline
\rowcolor{light_blue}
\textbf{技能項目} & \textbf{第1週} & \textbf{第2週} & \textbf{第3週} & \textbf{達成} & \textbf{備註} \\
\hline
心臟超音波基本操作 & $\square$ & $\square$ & $\square$ & $\square$ & 包含機器調整、探頭使用 \\
\hline
Apical 4-chamber view & $\square$ & $\square$ & $\square$ & $\square$ & 標準切面獲取 \\
\hline
Parasternal long axis view & $\square$ & $\square$ & $\square$ & $\square$ & 影像最佳化 \\
\hline
Parasternal short axis view & $\square$ & $\square$ & $\square$ & $\square$ & 多層面掃描 \\
\hline
Subcostal view & $\square$ & $\square$ & $\square$ & $\square$ & 下腔靜脈評估 \\
\hline
彩色都卜勒應用 & $\square$ & $\square$ & $\square$ & $\square$ & 血流評估技術 \\
\hline
週邊血管超音波 & $\square$ & $\square$ & $\square$ & $\square$ & 頸動脈、下肢血管 \\
\hline
DVT檢查技術 & $\square$ & $\square$ & $\square$ & $\square$ & 壓迫試驗 \\
\hline
PAOD評估 & $\square$ & $\square$ & $\square$ & $\square$ & ABI測量 \\
\hline
床邊緊急超音波 & $\square$ & $\square$ & $\square$ & $\square$ & RUSH protocol \\
\hline
\end{longtable}

\subsection{導管室技能檢核}

\begin{longtable}{|>{\raggedright\arraybackslash}p{3.5cm}|>{\raggedright\arraybackslash}p{1.5cm}|>{\raggedright\arraybackslash}p{1.5cm}|>{\raggedright\arraybackslash}p{1.5cm}|>{\raggedright\arraybackslash}p{1.5cm}|>{\raggedright\arraybackslash}p{4cm}|}
\hline
\rowcolor{light_green}
\textbf{技能項目} & \textbf{第1週} & \textbf{第2週} & \textbf{第3週} & \textbf{達成} & \textbf{備註} \\
\hline
無菌技術操作 & $\square$ & $\square$ & $\square$ & $\square$ & 刷手、穿無菌衣 \\
\hline
動脈穿刺觀摩 & $\square$ & $\square$ & $\square$ & $\square$ & 橈動脈、股動脈 \\
\hline
靜脈穿刺練習 & $\square$ & $\square$ & $\square$ & $\square$ & 中心靜脈穿刺 \\
\hline
導管操作基礎 & $\square$ & $\square$ & $\square$ & $\square$ & 導絲、導管選擇 \\
\hline
造影技術理解 & $\square$ & $\square$ & $\square$ & $\square$ & 角度選擇、對比劑 \\
\hline
血流動力學監測 & $\square$ & $\square$ & $\square$ & $\square$ & 壓力測量、心輸出量 \\
\hline
IVUS/OCT操作觀摩 & $\square$ & $\square$ & $\square$ & $\square$ & 血管內影像 \\
\hline
併發症識別 & $\square$ & $\square$ & $\square$ & $\square$ & 穿刺併發症、對比劑反應 \\
\hline
術後照護 & $\square$ & $\square$ & $\square$ & $\square$ & 止血、監測 \\
\hline
緊急處置參與 & $\square$ & $\square$ & $\square$ & $\square$ & STEMI處理流程 \\
\hline
\end{longtable}

\section{每日學習記錄表}

\subsection{學習日誌格式}
學生需每日完成以下學習記錄：

\begin{tcolorbox}[
    colback=emergency_blue!5!white,
    colframe=emergency_blue,
    title=\textbf{每日學習記錄範本},
    fonttitle=\bfseries,
    boxrule=0.5pt,
    arc=3pt
]
\textbf{日期：}\_\_\_\_\_\_\_\_\_\_ \textbf{星期：}\_\_\_\_\_\_

\textbf{今日主要活動：}
\begin{itemize}[nosep]
    \item 
    \item 
    \item 
\end{itemize}

\textbf{新學到的技能/知識：}
\begin{itemize}[nosep]
    \item 
    \item 
    \item 
\end{itemize}

\textbf{印象深刻的病例：}
\begin{itemize}[nosep]
    \item 病例描述：
    \item 學習重點：
    \item 反思思考：
\end{itemize}

\textbf{今日困難/疑問：}
\begin{itemize}[nosep]
    \item 
    \item 
\end{itemize}

\textbf{明日學習目標：}
\begin{itemize}[nosep]
    \item 
    \item 
\end{itemize}

\textbf{指導醫師回饋：}\_\_\_\_\_\_\_\_\_\_\_\_\_\_\_\_\_\_\_\_\_\_\_\_

\textbf{自我評分（1-5分）：}\_\_\_\_\_\_
\end{tcolorbox}

\section{專業發展規劃}

\subsection{心臟科專科訓練路徑介紹}
\begin{itemize}[nosep]
    \item \textbf{內科專科醫師訓練}：4年住院醫師訓練
    \item \textbf{心臟內科次專科}：3年心臟科fellow訓練
    \item \textbf{進階次專科選擇}：
    \begin{itemize}[nosep]
        \item 介入性心臟學（Interventional Cardiology）
        \item 心臟電生理學（Cardiac Electrophysiology）
        \item 心臟影像學（Cardiac Imaging）
        \item 心臟衰竭與移植（Heart Failure \& Transplant）
        \item 先天性心臟病（Adult Congenital Heart Disease）
    \end{itemize}
\end{itemize}

\subsection{研究機會介紹}
\begin{itemize}[nosep]
    \item 臨床研究參與機會
    \item 基礎科學研究合作
    \item 國際會議發表經驗
    \item 學術論文撰寫指導
    \item 研究經費申請經驗
\end{itemize}

\subsection{國際交流機會}
\begin{itemize}[nosep]
    \item 海外fellow訓練機會
    \item 國際醫學中心見習
    \item 國際會議參與
    \item 跨國研究合作項目
\end{itemize}

\section{最終評估和認證}

\subsection{綜合能力評估}
計畫結束時將進行綜合評估：

\begin{longtable}{|>{\raggedright\arraybackslash}p{4cm}|>{\raggedright\arraybackslash}p{6cm}|>{\raggedright\arraybackslash}p{4cm}|}
\hline
\rowcolor{light_orange}
\textbf{評估項目} & \textbf{評估內容} & \textbf{評估方式} \\
\hline
病例報告 & 
\begin{itemize}[nosep]
\item 選擇印象深刻病例
\item 完整病例分析
\item 文獻回顧整合
\item 學習心得分享
\end{itemize} & 
15分鐘口頭報告 + Q\&A \\
\hline
技能示範 & 
\begin{itemize}[nosep]
\item 超音波操作示範
\item 影像判讀能力
\item 臨床整合思維
\end{itemize} & 
實際操作評估 \\
\hline
知識測驗 & 
\begin{itemize}[nosep]
\item 心血管基礎知識
\item 臨床決策能力
\item 指引熟悉程度
\end{itemize} & 
筆試 + 口試 \\
\hline
360度評估 & 
\begin{itemize}[nosep]
\item 指導醫師評價
\item 同儕回饋
\item 護理師評價
\item 自我評估
\end{itemize} & 
多來源回饋表單 \\
\hline
\end{longtable}

\subsection{認證標準}
\begin{itemize}[nosep]
    \item \textbf{優秀（Excellent）}：總分90分以上
    \item \textbf{良好（Good）}：總分80-89分
    \item \textbf{合格（Satisfactory）}：總分70-79分
    \item \textbf{需改進（Needs Improvement）}：總分60-69分
    \item \textbf{不合格（Unsatisfactory）}：總分60分以下
\end{itemize}

\subsection{認證證書}
完成計畫並通過評估的學生將獲得：
\begin{itemize}[nosep]
    \item 心血管中心臨床見習完成證書
    \item 指導醫師推薦信
    \item 詳細學習成果報告
    \item 未來申請推薦支持
\end{itemize}

\begin{importantbox}
\textbf{計畫完成後支持：}
\begin{itemize}[nosep]
\item 提供住院醫師申請推薦信
\item 協助安排額外學習機會
\item 研究合作機會媒合
\item 專業發展諮詢服務
\item 持續mentor關係維持
\end{itemize}
\end{importantbox}

\section{聯絡資訊和緊急聯絡}

\subsection{主要聯絡人}
\begin{itemize}[nosep]
    \item \textbf{指導醫師}：洪國軒醫師、蕭喻中醫師
    \item \textbf{辦公室電話}：（03）5326151 分機 522009, 522010
    \item \textbf{行動電話}：0972-654046
    \item \textbf{電子郵件}：G11135@hch.gov.tw
\end{itemize}

\subsection{緊急聯絡}
\begin{itemize}[nosep]
    %\item \textbf{值班醫師}：（03）XXX-XXXX 分機XXXX
    \item \textbf{護理站}：（03）5326151 分機525022, 525002
    %\item \textbf{保全室}：（03）XXX-XXXX 分機XXXX
    \item \textbf{總機}：（03）5326151 分機9
\end{itemize}

\subsection{重要地點}
\begin{itemize}[nosep]
    \item \textbf{心臟內科門診}：醫療大樓1樓
    \item \textbf{心導管室}：醫療大樓2F
    \item \textbf{心血管病房}：醫療大樓2F樓2B ward
    \item \textbf{超音波室}：醫療大樓6樓
    \item \textbf{急診部}：醫療大樓1樓
\end{itemize}

\section{附錄}

\subsection{推薦閱讀清單}
\textbf{教科書：}
\begin{itemize}[nosep]
    \item Braunwald's Heart Disease: A Textbook of Cardiovascular Medicine
    \item Hurst's The Heart Manual of Cardiology
    \item The ESC Textbook of Cardiovascular Medicine
\end{itemize}

\textbf{臨床指引：}
\begin{itemize}[nosep]
    \item AHA/ACC Guideline for the Management of Patients with STEMI
    \item ESC Guidelines for acute and chronic coronary syndromes
    \item Taiwan Society of Cardiology Guidelines
\end{itemize}

\textbf{線上資源：}
\begin{itemize}[nosep]
    \item American College of Cardiology (ACC.org)
    \item European Society of Cardiology (ESCardio.org)
    \item Taiwan Society of Cardiology (TSOC.org.tw)
\end{itemize}

\subsection{學習APP推薦}
\begin{itemize}[nosep]
    \item \textbf{ECG學習}：ECG Guide by QxMD
    \item \textbf{心臟解剖}：3D Heart Anatomy
    \item \textbf{藥物查詢}：Epocrates, Micromedex
    \item \textbf{計算工具}：CardioSmart, AHA eCards
\end{itemize}

\note{
本計畫將依據實際臨床狀況彈性調整，確保學生獲得最佳學習體驗。
}

% References section
\newpage
\section{參考文獻}
\begin{thebibliography}{9}

\bibitem{aamc2019}
Association of American Medical Colleges. (2019). Core Entrustable Professional Activities for Entering Residency: Curriculum Developers' Guide. Washington, DC: AAMC.

\bibitem{mayo2020}
Mayo Clinic. (2020). Clinical Shadowing Program Guidelines. Mayo Clinic Medical Education.

\bibitem{cleveland2021}
Cleveland Clinic. (2021). Cardiovascular Medicine Training Program. Cleveland Clinic Academy.

\bibitem{hopkins2020}
Johns Hopkins Medicine. (2020). Clinical Skills and Simulation Center Training Modules. Johns Hopkins University School of Medicine.

\bibitem{acc2019}
American College of Cardiology. (2019). COCATS 4 Task Force: Training Statement for Cardiovascular Medicine. Journal of the American College of Cardiology, 73(17), 2202-2230.

\end{thebibliography}

\end{document}